\section{Introduction}
\label{sec:introduction}

Modern language model development increasingly relies on smaller proxy models and intermediate training checkpoints to guide decisions about data mixtures, data sampling strategies, architectures, tokenizers, and other training choices.
For the results of these smaller training runs to guide training at scale, they must lead to making the same decisions as if the run had been performed with a fully-trained larger model.
This predictability has been previously studied in the monolingual setting \citep{magnusson_datadecide_2025}, where small proxy models provide good signal about which design choices are best for larger ones. Indeed, inability to use proxy models or earlier checkpoints in the monolingual setting  is typically attributed to failures of the specific benchmark rather than lack of signal from the proxy runs
\citep{heineman_signal_2025}.

Multilingual language model development introduces a host of factors that could cause proxy run results to become noisy. %, such as cross-language transfer and interference and undertraining of languages with smaller portions of the token budget.
Languages vary in data availability, tokenization cost, and scaling behavior \citep{petrov_language_2023, ahia_all_2023,mueninghof-data-constrained}, while cross-lingual transfer further couples decisions across languages, meaning that adding data for one language can help or hurt others depending on model capacity and language similarity \citep{chang_when_2023, longpre_atlas_2026}. 
The token budget of a proxy run is also divided among its languages: while a total token budget may look large, languages receiving small fractions of the share can easily be undertrained. 
%: in the settings we study, half of each budget is English and the remainder is divided among up to 49 languages in proportion to corpus size, and a 350M-parameter proxy trained on 50 languages gives a language holding 0.5\% of the non-English half about 86M training tokens. 
As a result, a benchmark that provides a reliable training signal for one language or scale may be uninformative for another, making the use of proxy runs for multilingual language model development questionable.%\CL{whether an evaluation of a language is usable at a proxy point depends on the training tokens of that language, a quantity that model size alone does not determine}.

We see empirical evidence of this already: many multilingual language--benchmark pairs are known to exhibit little measurable signal at small scales, with performance often remaining near chance or floor levels.
In our own sweep, task accuracy exceeds the chance level in 4.0\% of Belebele's language--size cells over the five sizes from 175M to 1.7B non-embedding parameters, in 2.2\% for Global-MMLU, and in 79.4\% for HellaSwag. A comparison read from a score at chance selects a candidate by noise, potentially leading to suboptimal decisions at the full cost of all these proxy runs.

We organize this study as three questions, in the order a developer asks them.
First, at which training-token counts of a language does an evaluation produce a score distinguishable from its own noise (\S\ref{subsec:rq0})? For accuracy tasks, we apply the chance test above; for per-language bits per byte, which has no chance level, we test whether the score changes consistently with model size and with training tokens.
Second, among the evaluations that pass this test, at which proxy sizes and training fractions does the ranking of model variants at the proxy point agree with their ranking at the reference point, the fully trained largest model (\S\ref{subsec:rq2})?
Third, which statistics computed from the proxy runs alone predict that agreement, so that a developer can assess a proxy without training the reference model (\S\ref{subsec:rq3})?

To study these questions, we pretrain a controlled multilingual model ladder that varies model size, model architecture and language coverage.
We evaluate models throughout training on \nbenchmarks multilingual benchmarks, spanning different answer formats and scoring protocols, including answer-letter and full-choice evaluation with log-probability- and BPB-based scoring.
We find that \TODO{[fill in with finding per question.]} 

This work makes the following main contributions:

\begin{itemize}
    \item A controlled \textbf{multilingual scaling ladder} and open model release, spanning model sizes, training checkpoints, language mixtures, design interventions, and seed replicates. We release all \nmodels pretrained models and a total of \nckpts intermediate checkpoints, together with their evaluation results on \nbenchmarktasks benchmark-language tasks.
    \item \CL{\textbf{Per-language minimum token counts for multilingual evaluations}: a table reporting, per benchmark family and answer format, the training tokens of a language at which the evaluation passes the chance test, together with item reformulations that raise the passing language subsets at 1.7B from 11 to 86 of 105 for Belebele (\S\ref{subsec:rq0}).}
    % Replaced bullet kept for reference:
    % \item A methodological \textbf{framework for characterizing multilingual evaluations} for model-development decisions, measuring how predictably benchmarks scale across model size and training progress, and how early they recover the decisions made at larger scale.
    \item An \textbf{analysis of decision reliability}\CL{: decision accuracy of proxy points against the 1.7B reference along model size, training fraction, and training tokens of the language, with agreement between seed replicates as the attainable level (\S\ref{subsec:rq2}), and a comparison of} cheap surrogates for ranking agreement across multilingual regimes (\S\ref{subsec:rq3}).
\end{itemize}